\FloatBarrier
\section{Joint Discussion}
\label{sec:discussion}

The shared finding is a separation between nominal density, geometric
fidelity, and task utility. The HPC classification study demonstrates it
without a full-scene detector adapter: PU-GCN leads both Chamfer comparisons,
yet PU-Net alone improves Line~B classification and no method improves
Line~A. The lab detection study then shows an additional sensitivity to
patch construction, measured-point preservation, and detector input processing.
Together, the branches make a geometry-only account of upsampling quality
insufficient for the tested downstream tasks.

The experimental designs are complementary but not interchangeable.
ModelNet40 measures how well a separately trained PointNet++ learns each
representation. Frozen KITTI insertion measures compatibility with a
pretrained detector. KITTI adaptation then tests whether representation-specific
training can reduce that mismatch. Neither OA and AP nor the two Line~A/B
scores should be averaged into a universal measure of upsampling quality.
Their denominators, target tasks, and training interventions differ.

The tested geometric and learned methods target surface organisation,
local feature expansion, or generative reconstruction
\cite{huang2013ear,yu2018punet,qian2021pugcn,kim2023puedgeformer,zhang2025pdans}.
Those objectives provide useful geometric priors, but they do not directly
optimise the particular PointNet++, PointRCNN, or CenterPoint decisions
evaluated here. A point may reduce average surface distance while changing
a neighbourhood, occupying a previously empty voxel, or replacing a
measurement inside a finite downstream budget.

The reception and voxel probes support this interpretation at different
levels of the pipeline. PointNet++ discards much local candidate support
at high density; CenterPoint can activate or truncate many additional
voxels; PointRCNN resampling changes which points reach the network.
These are measured properties, but they do not uniquely decompose every
AP or OA deficit. In particular, Line~B loses detection performance even
without voxel-cap activation, and a full-input PointRCNN path also harms
the original baseline. Distribution compatibility extends beyond a single
hard input cap.

The object-level analyses strengthen the comparison without equating its
different units. On ModelNet40, PU-GCN's CD advantage holds for more
than 93\% of test objects in each line, yet it does not produce the
best classification result. The mismatch is therefore not explained
only by a few geometric outliers. Conversely, Line~B PU-Net's HD
advantage is widespread and accompanies its small classification gain,
but this co-occurrence does not establish HD as the cause of that gain.

The lab object evidence distinguishes a different set of alternatives.
The same historical inputs can lose Cars and recover a Cyclist in one
scene; more effective crop points can coexist with a lost box, while
a recovered PointRCNN object can have fewer effective crop points.
Thus neither an all-failure narrative nor a simple local-count rule
fits the retained observations. The loss/recovery totals supply the
population context for these deliberately selected examples.

These findings are compatible with the local grouping of PointNet++
and PointRCNN and the voxel-based representation of CenterPoint
\cite{qi2017pointnetpp,shi2019pointrcnn,yin2021centerpoint}, but the
visualisations do not directly measure learned feature changes.
The strongest supported interpretation concerns the complete input
and task pipeline. The relative contributions of geometry, sampling,
normalisation, and optimisation would require further controlled
interventions to separate.

The improvement after detector adaptation is consistent with a mismatch
between generated inputs and the original detector's learned representation.
Observation preservation helps separately by retaining the actual sensor
rows. Their combination gives the best completed KITTI results, especially
for CenterPoint. The extended schedules then tested whether the original
three-epoch stopping point was sufficient to explain the remaining loss.
All four observed-first conditions improved further and met the specified
Car-validation plateau rule, yet their baseline deficits persisted
(Section~\ref{sec:adaptation-evidence}). This strengthens the conclusion
within the tested schedules without excluding another optimisation or
input design. Because the baseline schedules were not extended to match,
and the validation set also selected checkpoints, the differences remain
conditional system comparisons rather than an isolated causal effect
of upsampling.

This distinction also prevents an inaccurate account of the later PU-GCN
work. The new contribution was the completed generation, controlled input
construction, detector adaptation, and full-validation evaluation. It was
not a newly trained PU-GCN model. Lab PU-EdgeFormer compatibility results
and flawed old PU-Net inputs likewise remain evidence about particular
implemented pipelines, rather than intrinsic rankings of their published
architectures.

The strongest evidence consists of complete paired dataset evaluations,
native-cardinality controls, and directly measured implementation effects.
The principal limits are one classifier training seed, test-based ModelNet40
checkpoint selection, batch-aggregated classifier metrics, provisional
mesh-reference training controls, unequal detector-training budgets,
and validation-based detector checkpoint selection. The plateau criterion
is specific to the chosen Car metric and the executed schedules. Diagnostic subsets, category
correlations, and selected visual cases refine the interpretation but do
not increase the number of independent primary experiments.

Within these limits, the practical implication is specific. An upsampling
pipeline for downstream perception must be assessed after the actual
input adapter and task model, with its observation-preservation rule,
reference geometry, and downstream training regime stated explicitly.
The work does not rule out useful upsampling. It identifies the conditions
under which the tested geometric gains failed to become task gains and
the partial recovery achieved by the completed adaptation procedure.
